\section{Note sur l'utilisation de l'intelligence artificielle}
Conformément aux directives académiques, l'auteur déclare avoir utilisé des outils 
d'intelligence artificielle générative (notamment Gemini 3.5 Flash et 3.1 Pro) durant la 
phase de rédaction de ce rapport intermédiaire.

L'utilisation de cet outil a été strictement encadrée et ciblée sur les points suivants :
\begin{itemize}
    \item \textbf{Assistance à la mise en page et structuration LaTeX :} Génération de structures de code
    (tableaux, indexation, glossaire) et résolution d'erreurs de compilation.
    \item \textbf{Aide à la formulation et synthèse :} Aide à la structuration logique du document, et 
    reformulation de notes de laboratoire brutes pour améliorer la clarté et la syntaxe de la rédaction.
    \item \textbf{Optimisation et génération de blocs de code :} Co-rédaction et ajustement de certaines portions de 
    code spécifiques (telles que les configurations de pipelines GStreamer ou des structures en C), ainsi que l'assistance 
    lors de la réalisation de scripts ou de code (C ou Python). Ces éléments ont été générés 
    systématiquement sur la base des propositions algorithmiques et des contraintes logiques dictées par l'auteur.
    \item \textbf{Soutien à la recherche documentaire :} Aide à l'analyse et à la navigation dans le 
    code source du pilote Morse Micro, notemment avec l'aide de l'outil Antigravity de google.
\end{itemize}

\textit{Note importante : L'intégralité de l'ingénierie système, de l'architecture globale, de la démarche itérative, de 
la conception des bancs de test matériels, des protocoles de mesure et des méthodologies de profilage n'a fait l'objet d'aucune 
génération par IA et provient exclusivement des travaux de l'auteur. De même, l'analyse des journaux d'exécution (logs) du noyau 
Linux a été réalisée et validée personnellement. Les requêtes (prompts) soumises à l'IA étaient systématiquement fondées sur ces 
recherches de bas niveau, des données brutes réelles et sur les notes de travail journalières issues des discussions avec Monsieur 
Mahboob Karimian (Chargé de R\&D, Institut IICT, HEIG-VD).}